\scanpage{27}
pour autant d'exister, et de définir des polygônes réguliers en ces car. exceptionelles, mais à un nb plus petit de cotés (divisant $n$). Ce serait \unsure{empoisonnant|emprisonnant} -- et \inserted{\text{finalement tout à fait}} artificiel -- de traîner tout au long du chemin la condition de fidélité, au lieu de se borner à l'examiner seulement aux moments \deleted{voulus} où cette considération est utile.

\textcircled{V} \deleted{\unsure{Abandon de $n$}} \uline{Cas des $\infty$-polygônes}

On peut aller encore plus loin, en notant que la donnée de (11) (a) revient à la donnée de
\[
\deleted{\unsure{a}}\quad
\left\{\begin{aligned}
&\sigma_E,\sigma'_E\in\mathrm{Aff}(E)\qquad\text{satisfaisant}\\
&\sigma_E^2={\sigma'_E}^2=\mathrm{id}_E,\quad
(\underbrace{\sigma'_E\sigma_E}_{u_E})^n=\mathrm{id}_E
\end{aligned}\right.\tag*{(14)}
\]
\deleted{satisfaisant} \deleted{(15) $s_0$} \deleted{Et la fidélité signifie que (étant don-} \deleted{\unsure{$s_0$}} \unsure{Etant|Étant} donné $s_0$, \deleted{la fidé} satisfaisant

\[
\deleted{(15)}\quad
\left\{\begin{aligned}
&\sigma_E s_0=s_0\\
&s_0,\ \deleted{\unsure{\sigma'_E\sigma_E s_0=u}}\,s_1=u_E s_0\text{ et }\inserted{s_2={}}u_E^2s_0\\
&\hspace{43mm}\text{sont non alignés}
\end{aligned}\right.\tag*{(13)}
\]
% archive sheet number: 26
